\section*{Appendix A. Complete Mathematical Formulations}\addcontentsline{toc}{section}{Appendix A}

\subsection*{Conditioning and Numerical Bounds}
Let $\kappa(W) = \sigma_1 / \sigma_\rho$ denote the matrix condition number. In \texttt{proxsparse\_pruner}, spectral truncation bounds intermediate condition ratios:
\[
\kappa_r(W_r) = \frac{\sigma_1}{\sigma_r}
\]

\subsection*{Gradient Dynamics of Decoupled Adapters}
During recovery adaptation, gradients flow according to:
\[
\frac{\partial \mathcal{L}}{\partial A} = \frac{\alpha}{r_{\text{lora}}} B^\top \left( \frac{\partial \mathcal{L}}{\partial Y} \right)^\top X, \quad \frac{\partial \mathcal{L}}{\partial B} = \frac{\alpha}{r_{\text{lora}}} \left( \frac{\partial \mathcal{L}}{\partial Y} \right) (X A^\top)
\]


\subsection*{Extended Sensitivity Proofs and Error Bounds}
\textbf{Proposition 3 (Lipschitz Continuity under Low-Rank Perturbation).} \textit{Let $f(X) = \sigma(X W^\top)$ be a neural network layer with $L_\sigma$-Lipschitz activation $\sigma$. Then the output perturbation under low-rank factorization is bounded by:}
\[
\|f(X) - f_r(X)\|_F \le L_\sigma \|X\|_2 \|W - W_r\|_F \le L_\sigma \|X\|_2 \sqrt{\sum_{i=r+1}^\rho \sigma_i^2}
\]
\textbf{Proof.} By the Lipschitz property of $\sigma$:
\[
\|\sigma(X W^\top) - \sigma(X W_r^\top)\|_F \le L_\sigma \|X (W - W_r)^\top\|_F \le L_\sigma \|X\|_2 \|W - W_r\|_F
\]
Applying the Eckart-Young truncation theorem completes the result. $\square$
